% А1.1 — Формулы сокращённого умножения; вынесение общего множителя и группировка
% Старая методичка: стр. 49 (старое задание 6)

\begin{tema}{А1.1}
\begin{pfTablica}[bez shapki]{rlrl}
\pfPodshapka{Квадраты и кубы суммы и разности}\\
\nomer1 & $(a+b)^2=a^2+2ab+b^2$ & \nomer3 & $(a+b)^3=a^3+3a^2b+3ab^2+b^3$\\
\nomer2 & $(a-b)^2=a^2-2ab+b^2$ & \nomer4 & $(a-b)^3=a^3-3a^2b+3ab^2-b^3$\\
\pfPodshapka{Разности степеней}\\
\nomer5 & $a^2-b^2=(a-b)(a+b)$ & \nomer6 & $a^3-b^3=(a-b)(a^2+ab+b^2)$\\
 & & \nomer7 & $a^3+b^3=(a+b)(a^2-ab+b^2)$\\
\pfPodshapka{Квадрат трёхчлена}\\
\nomer8 & \Block[l]{1-3}{$(a+b+c)^2=a^2+b^2+c^2+2ab+2bc+2ac$}\\
\end{pfTablica}

\opr{Вынесение общего множителя:} $ab+ac=a(b+c)$.
\qquad\opr{Группировка:} $ax+ay+bx+by=(a+b)(x+y)$.

\begin{primer}
$x^3+2x^2-x-2=x^2(x+2)-(x+2)=(x+2)(x^2-1)=(x+2)(x-1)(x+1)$
— группировка, затем формула (5).
\end{primer}
\end{tema}
